% !TEX root = ../main.tex
\chapter{Money and the State: International}\label{ch:intl}
\begin{paracol}{2}

\begin{sources}
\begin{tabularx}{\linewidth}{@{}l l L@{}}
\toprule
\textbf{Segment} & \textbf{Length} & \textbf{Content}\\
\midrule
\seg{14}{1}{L14.1 FT: Costs of Japan's Monetary Policy} & 2:49 & Zombie firms and a strong yen.\\
\seg{14}{2}{L14.2 Reading: Robert Mundell} & 10:27 & Mundell's 20th century in the language of the hierarchy: central banks avoiding the discipline of gold.\\
\seg{14}{3}{L14.3 Confrontation} & 11:48 & Act one: the Fed and the gold standard, 1931.\\
\seg{14}{4}{L14.4 Contradiction} & 14:08 & Act two: national management versus fixed rates; Keynes's bancor plan versus the IMF.\\
\seg{14}{5}{L14.5 The Dollar System} & 7:04 & Kindleberger: the US as a bank to the world; the end of Bretton Woods.\\
\seg{14}{6}{L14.6 Flexible exchange} & 8:28 & Act three: floating rates, inflation, inflation targeting, speculative FX markets.\\
\seg{14}{7}{L14.7 Global Financial Crisis} & 17:58 & European banks funding US mortgages in dollars; central bank swap lines; return of home bias.\\
\bottomrule
\end{tabularx}

\smallskip
\textbf{Files.} Transcripts \file{materials/transcripts/L14-P*.txt}; Mehrling's notes \mn{14}{1}.
\textbf{Reading.} R.~A.~Mundell, ``A Reconsideration of the Twentieth Century'', \emph{American Economic Review} 90(3), 2000, 327--340.
\end{sources}

\section*{Motivation}
Lecture~\ref{ch:chartalism} placed the dollar at the top of a hierarchy of currencies. This lecture tells the history of the international monetary system in
the 20th century, following Mundell's Nobel lecture, and translates it into the language of the course: discipline versus elasticity at the top of the
hierarchy \ts{14}{2}{3:12}. It ends with the global financial crisis, where the same tension reappears as dollar funding of non-US banks and
central bank swap lines \ts{14}{7}{0:03}.

% ------------------------------------------------------------------
\section{FT: the costs of Japan's monetary policy}\label{sec:intl-ft}

\begin{casestudy}[FT, 2012: John Plender, ``Japan counts zombie cost of ultra-loose monetary moves'']
About 20 years earlier Japan had a financial crisis; its responses included quantitative easing, the first experiment with this kind of unconventional
policy, and fiscal expansion. The economy never really got going again and has had deflation ever since \ts{14}{1}{0:07}.
Plender worries about the structural damage of long periods of low rates: uncompetitive firms do not go out of business, because rolling over their loans
costs them almost nothing, and by competing they keep the productive firms from making profits \ts{14}{1}{0:50}. The
discipline avoided for so long now comes in through the back door: a yen at 79--80 per dollar makes Japanese exports uncompetitive
\ts{14}{1}{1:53}.
\end{casestudy}

\begin{definition}[New concepts: zombie firm]
A firm that is not viable but survives because cheap credit lets it roll over its debts (the living dead of horror films: neither alive nor buried)
\ts{14}{1}{1:11}.
\end{definition}

% ------------------------------------------------------------------
\section{Mundell's reconsideration of the twentieth century}\label{sec:intl-mundell}

\paragraph{Still no global currency.} Mundell wrote the article in 1999, published in 2000 when he received the Nobel prize; it ends optimistically: after a bad
century we now have stable units, the euro, the dollar and the yen, and perhaps we are on our way to an international monetary system
\ts{14}{2}{0:03}. Twelve years later the yen and the euro still fluctuate a lot against the dollar and against each other, there is no global
currency, the euro (born in 1999) is in crisis, and after the financial crisis there has been some retreat from financial globalization towards ``national
capitalism'' \ts{14}{2}{0:44}.

\begin{coursenote}
On the FRED chart shown, the euro is quoted as dollars per euro (American convention) and the yen as yen per dollar (European convention), so a rise of the
first line and a fall of the second both mean appreciation against the dollar. With no movement between yen and euro the two lines would move in opposite
directions \ts{14}{2}{1:04}.
\end{coursenote}

\paragraph{The good old days of gold.} For Mundell the international gold standard is the ``good old days'' because the payment system disciplined central banks
and finance ministries, and, since everyone adjusted to the same discipline, the whole world was integrated and stable \ts{14}{2}{3:32}.
There were problems, notably a tendency to long-run deflation when the value of gold changed, but in retrospect the system looks good. It broke down in World
War~I \ts{14}{2}{4:55}.

\paragraph{Central banks avoiding discipline.} The 20th-century story is of countries, central banks and finance ministers who found the discipline too harsh and tried
to get out from under it \ts{14}{2}{5:36}. In the hierarchy gold $>$ currency $>$ deposits, the discipline sits in the link between gold and
currency: as the economy grows, currency expands but gold does not \ts{14}{2}{6:20}. Mundell keeps in mind two clean ways out:
\begin{enumerate}
  \item \textbf{revalue gold}: a higher price of gold is like more gold; but all countries must do it together, and since gold holdings differ it is a wealth
        transfer, which international politics could not agree on \ts{14}{2}{6:46};
  \item \textbf{create artificial gold}: increase the quantity instead of the value, as attempted with the SDR at the IMF, but too little too late
        \ts{14}{2}{7:47}.
\end{enumerate}
Neither was found, so there was ``a lot of bumbling around''. The main culprit is the Fed: a new central bank (1913), with no history, constantly making errors;
but the US is dominant, everyone pegs to the dollar, and what the Fed does the world does. Instead of a gold standard, a dollar standard with an immature Fed
\ts{14}{2}{8:27}. Then floating, until in 1999 Europe decided it no longer wanted floating rates among its currencies
\ts{14}{2}{9:32}.

\begin{nfigure}
\begin{tikzpicture}[font=\scriptsize]
  \draw[thick,->] (0,0) -- (14.5,0);
  \foreach \x/\y in {0/1900,4.8/1933--34,10.2/1971--72,14/1999}{\draw (\x,0.1) -- (\x,-0.1) node[below]{\y};}
  \fill[thmcol!20] (0,0.25) rectangle (4.6,1.35);
  \fill[fincol!20] (5.0,0.25) rectangle (10.0,1.35);
  \fill[defcol!20] (10.4,0.25) rectangle (14,1.35);
  \node[align=center] at (2.3,0.8) {Act 1: confrontation\\Fed vs gold standard};
  \node[align=center] at (7.5,0.8) {Act 2: contradiction\\national management vs fixed rates};
  \node[align=center] at (12.2,0.8) {Act 3: flexible rates\\learning from experience};
  \node[below,align=center] at (2.3,-0.5) {1913 Fed; 1914--18 war;\\1931 Fed raises rates};
  \node[below,align=center] at (7.5,-0.5) {1944 Bretton Woods; \$35/oz;\\1971 gold window closed};
  \node[below,align=center] at (12.2,-0.5) {1970s stagflation;\\inflation targeting; euro 1999};
\end{tikzpicture}
\caption{Mundell's 20th century in three acts, with the dates used in the lecture \ts{14}{3}{0:03}\ts{14}{4}{0:45}\ts{14}{6}{1:12}.}\label{fig:intl-acts}
\end{nfigure}

\begin{reading}[R.~A.~Mundell, A Reconsideration of the Twentieth Century (2000)]
Mundell's Nobel lecture (delivered December 1999). Read with the hierarchy of money in mind: every episode is a central bank, at the top of a national
hierarchy, facing an international standard above it \ts{14}{2}{9:52}.

Some facts from the article (\file{materials/readings/L14-Mundell.pdf}), checked against the scan:
\begin{itemize}
  \item Cassel had warned in 1925 and 1928 that restoring gold would bring a scarcity of gold and falling prices (p.~329).
  \item Wholesale prices fell from their 1929 peak to September 1931 by between 22\% (Germany) and 40.5\% (Japan); the United States figure is 29.5\% (p.~329).
  \item In October 1931 the Fed raised the rediscount rate in two steps from 1\textonehalf\ to 3\textonehalf\%; US wholesale prices fell 35\% between 1929 and 1933 (p.~330).
  \item Bank failures rose from about 500 a year in the 1920s to 1{,}350 in 1930, 2{,}293 in 1931 and 1{,}453 in 1932 (p.~330).
  \item In 1934 the dollar's gold value was cut by 40.94\%, raising the official gold price by 69.33\% to \$35 an ounce (p.~331).
\end{itemize}
Both quotations used below (p.~330 on ``pumping liquidity'', p.~331 on the counterfactual) are verbatim in the article.
\end{reading}

\begin{history}[Mundell]
Robert A.~Mundell (1932--2021), Canadian economist, received the 1999 Nobel memorial prize for his analysis of monetary and fiscal policy under different
exchange rate regimes and of optimum currency areas; he is often called an intellectual father of the euro.\par
\emph{References:} NobelPrize.org, ``Robert A.~Mundell, Facts'', 1999; R.~A.~Mundell, \emph{AER} 90(3), 2000.
\end{history}

\begin{exercise}[Two ways to more gold]
World gold reserves are fixed while nominal incomes must double. Show the two adjustments Mundell has in mind (price of gold, price level) and explain who gains
and who loses from each when gold holdings are unevenly distributed.
\end{exercise}

% ------------------------------------------------------------------
\section{Act one: confrontation (1900--1933)}\label{sec:intl-act1}

Mundell calls the first act the confrontation of the Fed and the gold standard; it really begins only after World War~I \ts{14}{3}{0:03}.
The Fed is created in 1913; almost immediately there is war, and every country goes off gold except the US. That was no feat: gold was flowing to the US to pay
for war material \ts{14}{3}{1:05}. So gold did not discipline the US during the war, and there was some wartime inflation; the discipline came back after the
war, with deflation in 1920--21, though not enough to return to pre-war price levels \ts{14}{3}{1:47}. One by one the mark,
the pound and the franc returned to gold, not all at the pre-war parities (the French devalued) \ts{14}{3}{2:28}.

\paragraph{Gold becomes scarce.} Mundell's point: each currency that re-links to gold adds to the demand for gold reserves, so gold becomes increasingly scarce,
undervalued \ts{14}{3}{3:13}. A first chance to revalue gold was missed. There are two ways to have ``more gold'': raise the price of gold, or keep
it and lower the price of everything else, which is deflation, bad for the economies and resisted by national central banks \ts{14}{3}{3:54}.

\paragraph{The 1931 mistake.} The key mistake, for Mundell, was the Fed's in 1931: to defend the gold standard it raised interest rates in the middle of a
deflation, from 1.5\% to 3.5\%, which in deflation are very high real rates \ts{14}{3}{4:40}. Gold flowed to the US, forcing the
other countries off gold, eventually the US too, and the depression became worldwide. Defending the parity of the dollar destroyed the gold standard, and at
home it sacrificed the US banking system, which collapsed; Roosevelt was elected and the banking legislation of the 1930s followed \ts{14}{3}{5:42}.

\begin{history}[The Fed in October 1931]
After Britain left gold (21 September 1931), the Federal Reserve Bank of New York raised its discount rate from 1.5\% to 2.5\% on 9 October and to 3.5\% on
16 October 1931 to stop the gold outflow, while US banks were failing in large numbers.\par
\emph{References:} M.~Friedman and A.~J.~Schwartz, \emph{A Monetary History of the United States, 1867--1960}, Princeton UP, 1963, ch.~7; B.~Eichengreen,
\emph{Golden Fetters}, Oxford UP, 1992, ch.~9.
\end{history}

\begin{intuition}[Two pars, one choice]
There are two par relationships: the mint par (dollar to gold) and the deposit par (deposits to currency). Defending the first, the Fed sacrificed the second:
people preferred currency to deposits, banks were run and closed. Alternatively it could have stopped paying gold and supported the banks as lender of last
resort. ``You can't do both'': lender of last resort to your banks, or defender of the dollar's value in world markets \ts{14}{3}{8:48}.
\end{intuition}

The notes quote Mundell (p.~330): ``Instead of pumping liquidity into the system, it chose to defend the gold standard'', compounding the secular deflationary tendency (every return to gold
raised the monetary demand for gold, so lowered gold prices of everything else; France, which devalued, got a fresh start) with explicitly deflationary policy. Falling farm prices meant loan
defaults, tight policy added liquidity pressure, and bank failures added a further deflationary impulse (on which the notes cite Irving Fisher, ``The Debt-Deflation Theory of Great Depressions'',
1933) \mn{14}{1}\mn{14}{2}.

\paragraph{A monetary theory of history.} On p.~331 Mundell writes that had the price of gold been raised in the late 1920s, or had the major central banks pursued
price stability instead of adhering to the gold standard, there would have been no Great Depression, no Nazi revolution and no World War~II
\ts{14}{3}{6:44}. Mehrling finds this overboard but dramatic: the point is that instability of the world polity and of the world
monetary system are tied together. As a student of politics he would say that both are symptoms of the same failure: we could not get together to do the right
thing \ts{14}{3}{7:45}.

\paragraph{Why did the Fed err?} A further story: Benjamin Strong, who ran the New York Fed and knew all the world's central bankers, had died, and so had his adviser
Allyn Young; the young Fed had no deep bench, and in the crisis the knee-jerk conservative reaction was to defend gold. Some say that had Strong lived the
depression might not have happened; ``we'll never know'' \ts{14}{3}{9:53}.

\begin{coursenote}
In class the date of Strong's death is given as ``1929 I think''. Benjamin Strong died on 16 October 1928, in New York (he had long suffered from
tuberculosis). Allyn Young died in London on 7 March 1929, of pneumonia during an influenza epidemic, as said. The biography mentioned is L.~V.~Chandler,
\emph{Benjamin Strong, Central Banker}, Brookings, 1958 \ts{14}{3}{10:14}.
\end{coursenote}

% ------------------------------------------------------------------
\section{Act two: contradiction (1934--1971)}\label{sec:intl-act2}

The second act is the confrontation of national macroeconomic management with a fixed exchange rate system, set up at Bretton Woods but starting already in
the late 1930s \ts{14}{4}{0:03}. Every finance minister and central bank wants policies good for its own country, and at the same
time commits to fixed rates with other countries: this cannot be squared unless someone absorbs all the contradiction, which would have been the US, and it
did not \ts{14}{4}{1:06}.

\paragraph{Keynes's bancor plan.} Before Bretton Woods (a resort in New Hampshire where the post-war plan was made), Keynes proposed a world central bank and a true
global currency, bancor \ts{14}{4}{1:48}. Deficit countries would borrow bancor, created out of thin air by the
international bank and credited to the surplus countries; so the quantity of international money expands with the imbalances. It is the classic banking
solution, a lender of last resort at the international level, just as inside the US a deficit agent ultimately borrows from a bank that borrows from the central
bank \ts{14}{4}{3:19}.

\begin{definition}[New concepts: bancor, asymmetry of adjustment]
\emph{Bancor}: Keynes's proposed international bank money (from French \emph{banque} and \emph{or}, ``bank gold''). \emph{Asymmetry of adjustment}: under gold
(and under the IMF) the deficit country must adjust, because it has to acquire the international reserve, while the surplus country need not: the survival
constraint at the level of nations \ts{14}{4}{6:07}.
\end{definition}

\begin{nfigure}
\begin{taccts}
\tacct[2.4cm]{Deficit country (UK)}{}{$+$Loan from int'l bank (bancor)}\quad
\tacct[2.4cm]{International bank}{$+$Loan to deficit country}{$+$Bancor deposit of surplus country}\quad
\tacct[2.4cm]{Surplus country (US)}{$+$Bancor deposit}{}
\end{taccts}
\caption{The bancor plan: imbalances expand the international bank's balance sheet \ts{14}{4}{3:19}.}\label{fig:intl-bancor}
\end{nfigure}

\begin{history}[Bretton Woods and Keynes]
The United Nations Monetary and Financial Conference met at Bretton Woods, New Hampshire, 1--22 July 1944; the US plan (Harry Dexter White) prevailed over
Keynes's Clearing Union. Keynes died on 21 April 1946, a month after the inaugural meeting of the IMF and World Bank governors at Savannah.\par
\emph{References:} B.~Steil, \emph{The Battle of Bretton Woods}, Princeton UP, 2013; R.~Skidelsky, \emph{John Maynard Keynes}, vol.~3, Macmillan, 2000.
\end{history}

\paragraph{Why the US said no.} The US asked: which of these agents is us? Clearly the surplus country, facing the UK, all of Europe, everyone else. The plan looked
like a mechanism forcing the US to lend to other countries, through an intermediary it did not control; true after the war, perhaps not forever
\ts{14}{4}{4:44}. Worse, to remove the asymmetry Keynes proposed charging interest on the surplus country's deposit as
well as on the loan, to push the US to spend its bancor on European goods: a forced loan at a negative rate \ts{14}{4}{6:07}.
It was an attempt to eliminate the survival constraint for the deficit countries, which needed US goods to rebuild \ts{14}{4}{7:29}.

\begin{coursenote}
Keynes's 1942--43 Clearing Union proposals did charge a fee on large \emph{credit} balances as well as on debit balances (at rates rising with the size of the
balance, of the order of 1--2\% a year); the 3\% in class is an illustrative number.\par
\emph{Reference:} ``Proposals for an International Clearing Union'', Cmd.~6437, HMSO, April 1943.
\end{coursenote}

\paragraph{What we got: the IMF.} The US plan. In stylised form (not the IMF's actual, more complicated accounting, seen in problem set~1), member countries
deposit gold and their own currencies and receive deposit-like claims, now called special drawing rights \ts{14}{4}{7:51}.
The Fund ends up with all kinds of currencies, the important one being the dollar \ts{14}{4}{10:44}. Deficit countries pay surplus countries in
SDRs: the quantity of international money does not change, it is only transferred, unlike bancor \ts{14}{4}{11:25}. When you hit zero you
borrow from the IMF, at a positive interest rate and with policy conditionality: exactly the asymmetry Keynes wanted to avoid \ts{14}{4}{12:06}.
Keynes died shortly after; perhaps he accepted the IMF thinking it could be fixed later (``the famous economist death theory of international finance'',
Mehrling's ``pet theory'') \ts{14}{4}{12:49}.

\begin{definition}[New concepts: special drawing right (SDR)]
An international reserve asset created by the IMF and held by member countries, which can use it to make payments to one another; ``paper gold'' of a fixed
quantity \ts{14}{4}{9:12}. (A \emph{drawing right} is a right to draw, to obtain other currencies, from the Fund.)
\end{definition}

\begin{coursenote}
The lecture's IMF balance sheet is a deliberate simplification. Historically, members paid their quota partly in gold and partly in their own currency and
could draw on the Fund up to limits; SDRs were created only later, by the First Amendment to the Articles (agreed 1967, in force 1969), with a first allocation
of about SDR 9.3 billion in 1970--72. They are allocated to members, not issued against deposits of gold.\par
\emph{Reference:} IMF, ``Special Drawing Rights (SDR)'' factsheet; M.~G.~de~Vries, \emph{The International Monetary Fund 1966--1971}, IMF, 1976.
\end{coursenote}

\paragraph{Discipline versus elasticity.} The IMF is a discipline system, bancor an elasticity system; maybe bancor would have been too elastic and the IMF too
disciplined. For a while the shortage of ultimate international reserves was mitigated by the expansion of dollars further down the hierarchy
\ts{14}{4}{13:10}.

\begin{exercise}[Bancor accounting]
The UK runs a deficit of 100 bancor with the US for three years. Write the balance sheets of the UK, the international bank and the US at the end, with
1\% interest charged on both debit and credit balances. Do the same under the stylised IMF, with initial SDR holdings of 150 each.
\end{exercise}


% ------------------------------------------------------------------
\section{The dollar system}\label{sec:intl-dollar}

Under Bretton Woods the hierarchy was: gold plus SDRs (a fixed quantity of paper gold) at the top; the dollar, pegged to gold at \$35 an ounce; the other
currencies (mark, pound, franc) pegged to the dollar \ts{14}{5}{0:03}. The extra reserves that allowed post-war growth came
neither from new gold, nor SDRs, nor revaluation, but from more dollars \ts{14}{5}{0:44}.

\begin{history}[Gold at \$35 and the closing of the gold window]
The Gold Reserve Act of January 1934 set the dollar price of gold at \$35 an ounce (from \$20.67). On 15 August 1971 President Nixon suspended the
convertibility of dollars held by foreign official institutions into gold. The view of the US as ``bank of the world'' was set out by E.~Despres,
C.~P.~Kindleberger and W.~S.~Salant, ``The Dollar and World Liquidity: A Minority View'', \emph{The Economist}, 5 February 1966.\par
\emph{References:} Federal Reserve History, ``Nixon Ends Convertibility of U.S. Dollars to Gold'' and ``Gold Reserve Act of 1934''.
\end{history}

\paragraph{The US as a bank (Kindleberger).} Mundell emphasizes trade deficits as the way dollars went abroad. Charles Kindleberger, Mundell's professor at MIT,
stressed a different mechanism, which needs no trade deficit at all: the US lends long term to the rest of the world (financing European reconstruction by
buying foreign bonds), while the rest of the world, which wants to grow, accumulates short-term dollar balances as monetary reserves
\ts{14}{5}{1:06}.

\begin{nfigure}
\begin{taccts}
\tacct[2.8cm]{United States (as a bank)}{Gold\\$+$Foreign bonds (long term)}{$+$Dollar deposits of the rest of the world (short term)}\qquad
\tacct[2.8cm]{Rest of the world}{$+$Dollar reserves}{$+$Bonds issued}
\end{taccts}
\caption{Kindleberger's view: the US as a bank to the world, borrowing short and lending long, with no net capital or trade flow \ts{14}{5}{2:42}.}\label{fig:intl-usbank}
\end{nfigure}

The US is a bank holding the world's international reserves, as long as the world accepts dollars and does not demand gold: a sort of bancor plan, but on the
US balance sheet; and the US earns the spread, because dollar reserves pay little or no interest \ts{14}{5}{3:03}.
The IMF refused to expand reserves at the top of the hierarchy; the US expanded them one step down \ts{14}{5}{4:04}.

\paragraph{The end of Bretton Woods.} Dollars expanded, gold did not, so gold and SDRs became undervalued and stress built up \ts{14}{5}{4:24}.
Again: revalue gold or create more SDRs. In 1967 there was an increase in SDRs, about \$10 billion, not enough. Mundell tells, for the first time in Mehrling's
reading, that Arthur Burns sounded out European central bankers about revaluing gold and suggested to Nixon to do it right after the 1968 election, which did
not happen \ts{14}{5}{4:44}. By 1971 there were too many dollars for the gold, there was a run on the dollar, and
instead of paying out the gold Nixon closed the gold window \ts{14}{5}{5:48}. Kindleberger argued all his life that it would have
been fine had foreigners accepted dollars as the reserve and not insisted on the ``fiction'' of gold \ts{14}{5}{6:15}.

\begin{exercise}[The US as a bank]
Starting from the T-accounts of figure~\ref{fig:intl-usbank}, show what happens when foreign central banks convert \$10 of their dollar balances into gold. Why is
the system a bank-run problem, and what are the options of the US?
\end{exercise}


% ------------------------------------------------------------------
\section{Act three: flexible exchange rates (1972--1999)}\label{sec:intl-act3}

The contradiction destroyed fixed rates; central bankers and finance ministers were now free to run their economies their own way, and the 1970s were not a good
time: stagflation, worldwide \ts{14}{6}{0:03}. Mundell calls 1972--1999 the experiment with flexible exchange rates and learning
from experience \ts{14}{6}{1:12}.

\begin{definition}[New concepts: stagflation]
Inflation and unemployment (stagnation) at the same time; the word joins \emph{stagnation} and \emph{inflation} \ts{14}{6}{0:23}.
\end{definition}

\paragraph{What was learned.} Flexible rates are very volatile, and without international discipline national policy tends to inflation: the greatest peacetime
inflation, because no central bank, least of all the Fed, faced any discipline \ts{14}{6}{1:40}. The stagnation came from
uncertainty: with exchange rates moving 100\% back and forth you cannot tell where your comparative advantage lies, and you invest in the wrong things
\ts{14}{6}{2:44}. Two consequences:
\begin{enumerate}
  \item \textbf{central bank self-discipline}: gold was too much discipline, but some is needed; central banks impose price discipline on themselves, inflation
        targeting \ts{14}{6}{3:27}. Friedman and the advocates of floating thought this would also stabilize exchange rates (PPP,
        lecture~\ref{ch:chartalism}); it did not: inflation differentials were small, yet exchange rates moved all over the place
        \ts{14}{6}{4:09}. Mundell wants central bank discipline \emph{and} fixed rates, central banks cooperating \ts{14}{6}{4:50};
  \item \textbf{private speculative markets}: firms selling abroad must hedge their exchange risk, so forward FX markets develop, very sophisticated; but the same
        markets allow speculation, such as the carry trade, which pushes currencies apart and then, when reversed, snaps them back
        \ts{14}{6}{4:50}. The dealers take risk and want to be paid for it, which distorts prices
        (next lecture) \ts{14}{6}{6:34}.
\end{enumerate}

\begin{definition}[New concepts: carry trade]
Borrowing in a low-interest currency and lending in a high-interest one (borrow yen, lend dollars) to pick up the interest differential, the \emph{carry},
bearing the exchange risk \ts{14}{6}{5:53}.
\end{definition}

\paragraph{Mundell's hope.} In 1999 Mundell foresaw three key currencies, each anchoring a large area (Asian currencies pegged to the yen, the dollar spreading
to Latin America, the euro taking over Europe), linked to each other into a unified world system. That is not what happened \ts{14}{6}{7:38}.

The notes add \mn{14}{4}\mn{14}{5}: a system of national currencies with floating rates is typical in wartime, when commerce is restricted, ``quite anomalous in peace time''. Friedman persuaded
himself that market prices beat administered ones and that domestic price stability would bring exchange stability through PPP; Mundell disagrees, valuing fixed rates as par clearing knits
together the states of the US. Inflation targeting did bring price stability, but not exchange stability: volatility of dollar, euro and yen continued. A lasting legacy of the period of disorder is
the currency futures market for hedging exchange risk (lecture~\ref{ch:futures}). And Mundell's three-currency picture is not symmetric: the world funding market is a dollar market, as the crisis showed.
In the 1960s, too, Mundell suggests that more SDRs, substituted for gold in reserves, would have lowered the demand for gold and taken the pressure off; one response to the recent crisis was indeed
a new issue of SDRs by the IMF \mn{14}{4}.

\begin{coursenote}
The IMF made a general SDR allocation of about SDR 161 billion (some \$250 billion) in August 2009, plus a special allocation in September 2009. The notes also mention, as an alternative
proposal of the 1940s, Benjamin Graham's commodity-reserve currency (\emph{World Commodities and World Currency}, 1944) \mn{14}{2}.\par
\emph{Reference:} IMF, ``Special Drawing Right (SDR)'' factsheet.
\end{coursenote}

% ------------------------------------------------------------------
\section{The global financial crisis}\label{sec:intl-gfc}

\paragraph{Dollar funding outside the US.} In the build-up to the crisis, big European banks (UBS, for example) bought US residential mortgage-backed securities and
funded them short term in the international dollar funding markets: shadow banking, though on balance sheet \ts{14}{7}{0:23}.
The money came, say, from money market mutual funds whose shares, a deposit substitute, were held by non-US residents, including central banks and countries that
prefer to keep dollars offshore, out of reach of the US legal system \ts{14}{7}{1:52}. Dollar borrowing and lending outside
the US, by institutions without access to the US banking system or to the Fed \ts{14}{7}{2:54}.

\paragraph{The run.} After Lehman and AIG, money funds that had lost on Lehman commercial paper wanted only Treasury bills and asked for their money back. The
European bank went to the ECB, which could lend euros but had no dollars: ``I'm not the Fed'' \ts{14}{7}{3:15}.

\paragraph{The swap-line operation.} The engineered solution \ts{14}{7}{5:20}:
\begin{enumerate}
  \item the Fed and the ECB swap deposits: the Fed creates dollars, the ECB creates euros (a liquidity swap);
  \item the ECB lends the dollars to the European bank, which rolls its funding and need not sell its mortgage securities;
  \item the dollars end up as a deposit of the US Treasury at the Fed, which issues additional T-bills to the money funds: the funds wanted T-bills, could not
        hold deposits at the Fed, and the Treasury acted as the intermediary.
\end{enumerate}

\begin{nfigure}
\begin{taccts}
\tacct[1.9cm]{European bank}{RMBS}{$-$MMF funding\\$+$Dollar loan from ECB}\quad
\tacct[1.9cm]{ECB}{$+$Dollar deposit at Fed\\$+$Dollar loan to bank}{$+$Euro deposit of Fed\\$-$Dollar deposit (lent)}\quad
\tacct[1.9cm]{Fed}{$+$Euro deposit at ECB}{$+$Dollar deposit of ECB\\(then Treasury deposit)}
\end{taccts}
\begin{taccts}
\tacct[1.9cm]{Treasury}{$+$Deposit at Fed}{$+$T-bills}\quad
\tacct[1.9cm]{Money market fund}{$-$Loan to European bank\\$+$T-bills}{Shares}
\end{taccts}
\caption{The central bank liquidity swap in the crisis, as drawn in class (changes only) \ts{14}{7}{4:37}.}\label{fig:intl-swap}
\end{nfigure}

\begin{definition}[New concepts: central bank liquidity swap line]
An agreement by which two central banks exchange their own currencies, to be swapped back later at the same exchange rate; the receiving central bank lends the
foreign currency on to its own banks \ts{14}{7}{6:27}.
\end{definition}

\begin{coursenote}
The Treasury deposit is the Supplementary Financing Program, announced in September 2008: the Treasury issued extra bills and left the proceeds on deposit at the
Fed. The Fed's swap lines peaked at about \$580 billion in December 2008 (``600 billion'' in class).\par
\emph{References:} Federal Reserve, H.4.1 releases 2008; Federal Reserve History, ``Central Bank Liquidity Swaps''.
\end{coursenote}

\paragraph{Temporary cooperation, then repatriation.} This is the central bank cooperation Mundell asks for, but only a holding action, to keep the shadow banks from
dumping their securities and collapsing the market; eventually the Fed bought mortgage securities itself \ts{14}{7}{7:54}. The
permanent solution was repatriation: bringing the securities back home, undoing the internationalization of finance; a return of \emph{home bias} in
investment, because lender-of-last-resort support is national \ts{14}{7}{8:56}. In the notes: the money fund still funds the shadow bank's MBS, but through a chain, money fund to
Treasury to Fed to ECB to shadow bank; afterwards the construction was dismantled by taking the MBS back to the US, where dollar funding was easier to cobble together: ``it is no accident
that there is now over a trillion MBS on the balance sheet of the Fed'' \mn{14}{5}\mn{14}{6}. The dollar is a national currency,
the liability of the Fed; there is no global currency \ts{14}{7}{9:58}.

\begin{definition}[New concepts: home bias]
The tendency of investors to hold domestic assets more than international diversification would suggest \ts{14}{7}{9:38}.
\end{definition}

The reversal of globalization worries Mehrling: we saw it at the end of the 1920s, and it is a headwind against which national policies do not work; countries
try fiscal or monetary stimulus and only become more indebted and vulnerable \ts{14}{7}{10:40}. Meanwhile the system is held together
by international central bank cooperation and by speculative trading \ts{14}{7}{11:42}.

\paragraph{Is the Fed taking risk?} A student asks; Congress asked Bernanke the same about the roughly \$600 billion of swaps \ts{14}{7}{12:03}.
Two answers:
\begin{enumerate}
  \item the swap is reversed at the same amounts: the Fed gets back its dollars, the ECB its euros. Each books in its own currency, so whatever the exchange rate
        does nobody books a loss; commercially, with both parties keeping their books in the same currency, one side would show a loss
        \ts{14}{7}{13:07};
  \item it is international lender-of-last-resort action: with private funding gone, the central banks had no real choice; they were doing their job (``try to
        explain this to Congress'') \ts{14}{7}{14:49}.
\end{enumerate}

\paragraph{Standing swap lines.} The big central banks (Japan, ECB, Switzerland, Bank of England, Fed) have agreed swap lines in all currencies, of unlimited amount:
they are aware of the history and of the stakes, that monetary instability means economic and then social instability \ts{14}{7}{15:50}.
Maybe they overestimate their influence: without political will, a lender of last resort cannot prevent social instability; wars happen when people cannot
figure out what to do next, and in war the central bank becomes the financing agent of the government. ``I'd rather not go down that route''
\ts{14}{7}{17:12}.

\begin{coursenote}
The coordinated action of 30 November 2011 involved six central banks (the five named plus the Bank of Canada). The swap lines became standing, unlimited
arrangements among these six only on 31 October 2013, after this lecture.\par
\emph{Reference:} Federal Reserve press releases of 30 November 2011 and 31 October 2013.
\end{coursenote}

% ------------------------------------------------------------------
\flushside

\begin{exercise}[Swap gains and losses]
The Fed swaps \$100 for euros at \$1.25 per euro; three months later the euro trades at \$1.40. Compute the gain or loss of each central bank if both kept their
books in dollars, and then if each kept its books in its own currency.
\end{exercise}
\section*{Key formulas and ideas}
\begin{keyformulas}
\begin{tabularx}{\linewidth}{@{}l L@{}}
Mundell's thesis & integration and stability of the world economy depend on the discipline of an international standard\\
Escapes from gold & revalue gold (price) or create paper gold (quantity); neither coordinated in time\\
1931 & defending the mint par sacrificed the deposit par: LLR to banks or defend the currency, not both\\
Bancor vs IMF & elasticity (reserves grow with imbalances) vs discipline (fixed SDRs, deficit country adjusts)\\
Dollar system & US as a bank: short-term dollar liabilities, long-term foreign assets, gold reserve\\
Floating & inflation, volatility; answers: inflation targeting, private hedging/speculative FX markets\\
Crisis & offshore dollar funding; swap lines; repatriation and home bias; no global currency\\
\end{tabularx}
\end{keyformulas}
